% Template for CRC journal article

%-----------------------------------------------------------------------------------

%-----------------------------------------------------------------------------------

%%%%%%%%%%%%%%%%%%%%%%%%%%%%%%%%%%%%%%%%%%%%%%
%%%%%%%%%%%%%%%%%%%%%%%%%%%%%%%%%%%%%%%%%%%%%%
%%                                          %%
%% Important note on usage                  %%
%% -----------------------                  %%
%% This file must be compiled with PDFLaTeX %%
%% Using standard LaTeX will not work!      %%
%%                                          %%
%%%%%%%%%%%%%%%%%%%%%%%%%%%%%%%%%%%%%%%%%%%%%%
%%%%%%%%%%%%%%%%%%%%%%%%%%%%%%%%%%%%%%%%%%%%%%

%% The '5p' and 'times' class options of elsarticle are used
\documentclass[5p,times,authoryear]{sciarticle}

%% The `crc' package must be called to make the CRC functionality available
%% crc_RIAI es el paquete ecrc para la revista RIAI
\usepackage{crc_SCI}

%% The crc package defines commands needed for running heads and logos.
%% For running heads, you can set the journal name, the volume, the starting page and the authors

%%%%%%%%%%%%%%%%%%%%%%%%%%%%%%%%% Añadido por Secretaría RIAI
\usepackage[spanish]{babel}     % Idioma
\addto\captionsspanish{%
\def\tablename{Tabla}%
}
\usepackage[utf8]{inputenc}     % Lengua latina
\usepackage{amsmath}            % Para las referencias a ecuaciones con \eqref
\usepackage{epstopdf}           % Para poder insertar figuras .eps al compilar con PDFLATEX
\usepackage{flushend}           % Para igualar las columnas de la última página
%\usepackage{hyperref}          % Para hipervínculos dentro del PDF
%%%%%%%%%%%%%%%%%%%%%%%%%%%%%%%%%%%%%%%%%%%%%%%%%%%%%%%


%% set the volume if you know. Otherwise `00'
\volume{XVIII}

%% set the starting page if not 1
\firstpage{1}

%% Give the name of the journal
\journalname{Simposio CEA de Control Inteligente.}

%% Give the author list to appear in the running head
%% Example \runauth{C.V. Radhakrishnan et al.}
\runauth{Apellido primer autor, Inicial. et al. }

%% The choice of journal logo is determined by the \jid and \jnltitlelogo commands.
%% A user-supplied logo with the name <\jid>logo.pdf will be inserted if present.
%% e.g. if \jid{yspmi} the system will look for a file yspmilogo.pdf
%% Otherwise the content of \jnltitlelogo will be set between horizontal lines as a default logo

%% Give the abbreviation of the Journal. Contast the Publisher if in doubt what this is.
\jid{SCI}

%% Give a short journal name for the dummy logo (if needed)
\jnltitlelogo{}

%% Hereafter the template follows `elsarticle'.
%% For more details see the existing template files elsarticle-template-harv.tex and elsarticle-template-num.tex.

%% The conventions of elsarticle-template-num.tex should be followed (included below)
%% If using BibTeX, use the style file elsarticle-num.bst

%% End of ecrc-specific commands
%%%%%%%%%%%%%%%%%%%%%%%%%%%%%%%%%%%%%%%%%%%%%%%%%%%%%%%%%%%%%%%%%%%%%%%%%%

%% The amssymb package provides various useful mathematical symbols
\usepackage{amssymb}
%% The amsthm package provides extended theorem environments
%% \usepackage{amsthm}

%% The lineno packages adds line numbers. Start line numbering with
%% \begin{linenumbers}, end it with \end{linenumbers}. Or switch it on
%% for the whole article with \linenumbers after \end{frontmatter}.
%% \usepackage{lineno}

%% natbib.sty is loaded by default. However, natbib options can be
%% provided with \biboptions{...} command. Following options are
%% valid:

%%   round  -  round parentheses are used (default)
%%   square -  square brackets are used   [option]
%%   curly  -  curly braces are used      {option}
%%   angle  -  angle brackets are used    <option>
%%   semicolon  -  multiple citations separated by semi-colon
%%   colon  - same as semicolon, an earlier confusion
%%   comma  -  separated by comma
%%   numbers-  selects numerical citations
%%   super  -  numerical citations as superscripts
%%   sort   -  sorts multiple citations according to order in ref. list
%%   sort&compress   -  like sort, but also compresses numerical citations
%%   compress - compresses without sorting
%%
%% \biboptions{comma,round}

% \biboptions{}

% if you have landscape tables
\usepackage[figuresright]{rotating}

% put your own definitions here:
%   \newcommand{\cZ}{\cal{Z}}
%   \newtheorem{def}{Definition}[section]
%   ...

% add words to TeX's hyphenation exception list
%\hyphenation{author another created financial paper re-commend-ed Post-Script}

% para poder introducir varias figuras que ocupen el ancho de las dos columnas.
\usepackage{subfigure}
\usepackage{float} % JVS
\usepackage{multicol} % JVS
\usepackage[most]{tcolorbox}


% declarations for front matter

\begin{document}

\begin{frontmatter}

%% Title, authors and addresses
%% use the tnoteref command within \title for footnotes;
%% use the tnotetext command for the associated footnote;
%% use the fnref command within \author or \address for footnotes;
%% use the fntext command for the associated footnote;
%% use the corref command within \author for corresponding author footnotes;
%% use the cortext command for the associated footnote;
%% use the ead command for the email address,
%% and the form \ead[url] for the home page:
%% \title{Title\tnoteref{label1}}
%% \tnotetext[label1]{}
%% \author{Name\corref{cor1}\fnref{label2}}
%% \ead{email address}
%% \ead[url]{home page}
%% \fntext[label2]{}
%% \cortext[cor1]{}
%% \address{Address\fnref{label3}}
%% \fntext[label3]{}
\dochead{XVIII Simposio CEA de Control Inteligente}
\datesite{28-30 de junio de 2023, Valencia}
%% Use \dochead if there is an article header, e.g. \dochead{Short communication}

% Escriba el título del artículo.
\title{Preparación de artículos para la revista RIAI:\\
no escriba mayúsculas en las palabras del título}


%% use optional labels to link authors explicitly to addresses:
%% \author[label1,label2]{<author name>}
%% \address[label1]{<address>}
%% \address[label2]{<address>}

\author[First]{Apellido1, Inicial1.\corref{cor1}}

\author[Second]{Apellido2, Inicial2.}

\author[Third]{Apellido3, Inicial3.}

% Escriba el email del autor para correspondencia.
\cortext[cor1]{Autor para correspondencia: autor1@ceautomatica.es
\\
Attribution-NonCommercial-ShareAlike 4.0 International (CC BY-NC-SA 4.0)
}

\address[First]{Comité Español de Automática, Parc Tecnologic de Barcelona, Edifici U, C/ Llorens i Artigas, 4-6, 08028 Barcelona, España. }
\address[Second]{Departamento de Automática, Ingeniería Electrónica e Informática, Universidad Politécnica de Madrid,  C/ José Gutiérrez Abascal, nº2, 28006, Madrid,  España.}
\address[Third]{Departamento de Ingeniería de Sistemas y Automática,  Universitat Politécnica de Valencia, Camino de Vera, nº14, 46022, Valencia, España.}

% Escriba el nombre de los autores y el título de su artículo en inglés para que aparezca correctamente el cuadro "Como citar:"
\begin{comocitar}
\begin{center}
\begin{tcolorbox}[frame hidden, width=\linewidth*2, colframe=white, colback=black!8!white]
\tocitearticle{Apellido1, Inicial1., Apellido2, Inicial2., Apellido3, Inicial3. 2020. Paper Title in English. Revista Iberoamericana de Automática e Informática Industrial 00, 1-5. https://doi.org/10.4995/riai.2020.7133}
\end{tcolorbox}
\end{center}
\end{comocitar}

\begin{abstract}
%% Text of abstract
Estas instrucciones constituyen una guía para la preparación de artículos para la revista RIAI. Utilice este documento como un conjunto de instrucciones. Por favor, utilice este documento como una ''plantilla'' para preparar su manuscrito. Para las directrices de envío, siga las instrucciones del sistema de envío de artículos de la página web de la revista. El título no debe tener más de 10 palabras. El número de autores no debe ser mayor de seis. El resumen no debe ocupar más de 160 palabras (aproximadamente de 8 a 10 líneas). Sustituya Apellido1, Inicial1. por el apellido del primer autor y la inicial de su nombre, Apellido2, Inicial2 por el del segundo, etc. Recuerde que el resumen no puede superar la primera página. En el encabezado de la página 2 y siguientes debe aparecer el apellido del primer autor y su inicial seguido por ''et al.''. No modifique el DOI y el número y páginas de la revista, ni en el encabezado, ni en el apartado To cite this article, cuando su artículo sea publicado la secretaría técnica de la revista será la encargada de hacerlo.
\end{abstract}

\begin{keyword}
Palabra 1 \sep Palabra 2 \sep 5-10 Palabras clave (tomadas de la lista del sitio web de IFAC y traducidas al castellano).
\end{keyword}

% Escriba el título en inglés.
\begin{englishtitle}
\noindent \textbf{Paper Title in English, Bold Style}
\end{englishtitle}

%% Escriba el resumen en inglés.
\begin{abstractIng}
These instructions constitute a guide for the preparation of articles for the RIAI magazine. Use this document as a set of instructions. Please use this document as a "template" to prepare your manuscript. For shipping guidelines, follow the instructions of the article submission system on the magazine's website. The title should not have more than 10 words. The number of authors should not be greater than six. The abstract should not occupy more than 160 words (approximately 8 to 10 lines). Replace Surname1, Initial1. by the last name of the first author and the initial of his name, Surname2, Initial2 by that of the second, etc. Remember that the summary cannot exceed the first page. In the header of page 2 and following should appear the last name of the first author and his initial followed by ''et al.''. Do not modify the DOI and the number and pages of the journal, neither in the heading, nor in the To cite this article section, when your article is published the technical secretariat of the journal will be in charge of doing so.
\end{abstractIng}

%% Escriba las palabras clave en inglés de la forma siguiente: palabra1 \sep palabra2 \sep palabra3
\begin{keywordIng}
Keyword 1 \sep Keyword 2 \sep 5-10 Keywords in English (taken from the keyword list at the IFAC website).
\end{keywordIng}


\end{frontmatter}

\begin{multicols}{2}   % JVS

%%
%% Start line numbering here if you want
%%
% \linenumbers

%% main text
\section{Introducción}
Estas instrucciones constituyen una guía para la preparación de
artículos para la revista RIAI. Utilice este documento como un
conjunto de instrucciones. Puede usar este documento como
una ``plantilla'' para preparar su manuscrito en Latex. Para las directrices
de envío, siga las instrucciones del sistema de envío de artículos
de la página web de la revista.
\emph {No cambie el tamaño de las fuentes o espaciado de línea para introducir más texto en un número limitado de páginas}.  Utilice \emph{cursiva} para enfatizar; no subraye.
Estas instrucciones constituyen una guía para la preparación de
artículos para la revista RIAI. Utilice este documento como un
conjunto de instrucciones. Puede usar este documento como
una ``plantilla'' para preparar su manuscrito en Latex.
\emph {No cambie el tamaño de las fuentes o espaciado de línea para introducir más texto en un número limitado de páginas}.  Utilice \emph{cursiva} para enfatizar; no subraye.
Es conveniente que lea y compruebe cada uno de los items que aparece en el fichero \textit{LEEME\_IMPORTANTE\_checklist.pdf} que incluye el lote de estilo descargado del sitio web de la revista.

\subsection{Una subsección}
Bifurcación: Trazado del máximo local de $x$ con una disminución de amortiguamiento $a$ (Figura~\ref{fig1}).

Las siguientes secciones incluyen instrucciones que deben seguir los autores sobre el procedimiento de envío de trabajos a la revista RIAI. Incluyen el procedimiento, la forma de introducir tablas y figuras y aspectos sobre las fases de revisión. También se establece el formato de las referencias bibliográficas.

RIAI no realizará ninguna operación de formateado final a su artículo. Su documento debe estar ``listo para filmar''. El número límite de hojas para el documento es de doce.

\begin{figure}[H]
\centering
  % El fichero es un eps y se convierte automáticamente a pdf con eps2pdf package
  \includegraphics[width=8cm]{Figuras/figuraeps}\\
  \caption{Título de la Figura 1. Recuerde que los títulos de figura que ocupan una sola línea van centrados. La figura es un fichero eps y gracias al paquete epstopdf se convierte automáticamente a pdf.}\label{fig1}
\end{figure}

\section{Procedimiento para el envío de artículos}
Recuerde que RIAI está considerado como un \emph{Camera Ready Copy Journal (CRC)}. Esto implica que los autores son responsables de aplicar el formato correspondiente a sus contribuciones. Desde la secretaría de la revista no se ejecutará ninguna acción de formateo a los artículos. A continuación vemos unas subsecciones.
\subsection{Fase de revisión}
Por favor, use este documento como una ``plantilla'' para preparar su documento. Para las directrices de envío, siga las instrucciones del sistema de envío de artículos.

Dado que el límite de páginas es de doce, es mejor preparar el envío inicial en el formato listo para filmar, de tal manera que tenga una buena estimación de la longitud de hojas. Adicionalmente, el esfuerzo requerido para el envío de la versión final será, de esta manera, mínimo.
\subsection{Fase final}
Se supone que los autores tendrán en cuenta rigurosamente los márgenes. En caso de no ser así se le pedirá que reenvíe el documento para que así lo cumpla, retrasando de esta manera la preparación de los contenidos de la revista \citep{Abl:45}, \citep{Abl:56}.

\subsection{Inserción de tablas}

La tabla ocupa el ancho de la columna porque el entorno \emph{tabular} lleva el asterisco (ver tabla \ref{extremos45}). Se puede usar \emph{table}* para confeccionar una tabla que se expanda sobre la dos columnas del texto (ver tabla \ref{tabladeseables}). Y por supuesto combinar ambos efectos \citep{Heritage:92}, \citep{ChaRou:66}.

\begin{table}[H]
  \caption{Preferencias para el diseño de un controlador}
   \label{extremos45}
  \begin{tabular*}{\hsize}{lrrrrr}
\hline
    & $g_i^1$ & $g_i^2$ & $g_i^3$ & $g_i^4$ & $g_i^5$ \\
    \hline
$Re(\lambda)_{max}$ & -0.01  & -0.005 & -0.001 & -0.0005 & -0.0001 \\
$u_{max}$& 0.85 & 0.90 & 1 & 1.5 & 2  \\
$t_{est}^{max}$& 14 & 16 & 18 & 21 & 25 \\
$noise_{max}$& 0.5 & 0.9 & 1.2 & 1.4 & 1.5  \\
$u_{nom}$& 0.5 & 0.7 & 1  & 1.5 & 2  \\
$t_{est}^{nom}$& 10 & 11 & 12 & 14 & 15 \\
\hline
  \end{tabular*}
\end{table}

Es muy importante mantener los márgenes. Son necesarios para poner información de la revista y los números de página.

En esta plantilla no se pueden usar \emph{figure} y \emph{table} con las opciones tradicionales: htbp. Solo se puede usar la opción "H". Si no se usa esta opción la figura o tabla no se incluirá en el documento.

\subsection{Figuras y creación del archivo PDF}

Todas las figuras deben estar incrustadas en el documento. Cuando incluya una imagen, asegúrese de insertar la imagen real en lugar de un enlace a su computador local. En la medida de lo posible, utilice las herramientas de conversión a PDF estándares Adobe Acrobat o Ghostscript que dan los mejores resultados. \emph{Es importante que todas las fuentes estén incrustadas en el PDF resultante}.

Al compilar utilizando PDFLatex, se pueden insertar figuras en jpg (Figura \ref{fig2}) o pdf (Figura \ref{fig3}). Si tiene figuras en formato eps conviértelas a pdf previamente o bien haga uso del paquete epstopdf.

Recuerde que esta plantilla solo usa la opción "H'' y que de otro modo la figura o tabla no se incluirá en el documento.

\section{Unidades}

Use el Sistema Internacional como unidades primarias. Se pueden usar
otras unidades como unidades secundarias (entre paréntesis). Esto se
aplica a artículos sobre almacenamiento de datos. Por ejemplo,
escriba ``$15 Gb/cm^2$'' ($100 Gb/in^2$). Se considera una excepción
cuando las unidades inglesas se usan como identificadores
comerciales, como unidad de disco de 3.5 pulgadas. Evite mezclar
unidades del Sistema Internacional con el Sistema Cegesimal, tales
como corriente en amperios y campo magnético en  oersteds. Esto a
menudo lleva a confusión porque las ecuaciones no son
dimensionalmente equiparables. Si debe usar unidades mezcladas,
especifique claramente las unidades para cada cantidad  en la
ecuación \citep{SanGra:99}.
\begin{figure}[H]
\centering
  % Se pueden incluir figuras jpeg al compilar con PDFLatex
  \includegraphics[width=3cm]{Figuras/figurajpeg}\\
  \caption{Título de la Figura 2. Recuerde que los títulos de figura que ocupan una sola línea van centrados.}\label{fig2}
\end{figure}

La unidad en el Sistema Internacional para la fuerza del campo
magnético H es A/m. Sin embargo, si desea utilizar unidades de $T$,
o bien refiérase a densidad de flujo magnético $B$ o fuerza del
campo magnético simbolizado como $\mu_0 H$. Utilice el punto
centrado para separar unidades compuestas,  es decir, $A\cdot m^2$.

\begin{figure*}
\centering
\subfigure[Título Subfigura 1]{
   \includegraphics[scale =0.5] {Figuras/figuraeps}
   \label{subfig1}
 }
 \subfigure[Título Subfigura 2]{
   \includegraphics[scale =0.35] {Figuras/figurajpeg}
   \label{subfig2}
 }
\label{mifigura}
\caption{Título de una sola linea al centro.}
\end{figure*}


\section{Consejos útiles}

\subsection{Figuras y tablas}

Cuando en el texto haga referencia a una figura o una tabla, recuerde que la palabra tabla o figura debe llevar mayúscula. Por ejemplo: Como se puede observar en la Figura \ref{mifigura}...; o (Tabla \ref{extremos45}). Recuerde que los títulos de figura que ocupan más de una línea van justificados. Si el título ocupa una sola línea va centrado, si el título ocupa más de una línea va justificado.
Las etiquetas de los ejes de las figuras son a menudo fuentes de
confusión. Utilice palabras en lugar de símbolos. Como ejemplo,
escriba la cantidad ``Magnetización,'' o ``Magnetización M,'' no
sólo ``M.'' Ponga las unidades entre paréntesis. No etiquete los
ejes únicamente con unidades. Por ejemplo,
escriba ``Magnetización (A/m)'' o ``Magnetización (A $\cdot$
m$^{-1}$),'' no sólo ``A/m'' No etiquete los ejes con una relación
de cantidades y unidades. Por ejemplo, escriba ``Temperatura (K),''
no ``Temperatura/K.''

Los multiplicadores pueden ser especialmente fuente de confusión.
Escriba ``Magnetización (kA/m)'' o ``Magnetización ($10^3$ A/m).'' No
escriba ``Magnetización (A/m) $\times$ 1000'' porque el lector no
sabría si la etiqueta es 16000 A/m o
0.016 A/m. Las etiquetas de las figuras deben ser legibles,
aproximadamente de 8 a 12 puntos.

Con el entorno \emph{figure*} se puede conseguir que una figura ocupe las dos columnas (Figura \ref{mifigura}). Con el paquete \emph{subfigure} conseguimos una figura completa a partir de varios ficheros (como las subfiguras \ref{subfig1} y \ref{subfig2}).

\begin{figure}[H]
\centering
  % Se pueden incluir figuras pdf al compilar con PDFLatex
  \includegraphics[width=8cm]{Figuras/figurapdf}\\
  \caption{Título linea simple al centro.}\label{fig3}
\end{figure}

De igual modo, el entorno \emph{table*} permite que una tabla ocupe las dos columnas (Tabla \ref{tabladeseables}).

En ambos entornos a dos columnas no se pueden usar ninguna de las opciones htbp de figure y table tradicionales y hay que dejar esta opción en blanco (por defecto usará "H''), pues de otro modo la figura o tabla no se incluirá en el documento.

\subsection{Referencias} % ATENCION: usar \citep en lugar de \cite para ajustarse al formato de RIAI.

La lista de referencias debe ser ordenada alfabéticamente
de acuerdo con el primer autor, con las siguientes líneas justificadas con la sangría correspondiente.
Si existen diferentes publicaciones del mismo autor(es), éstas deberán ser listadas en el orden del
año de publicación. Si hay más de un artículo del mismo autor en la misma fecha,
etiquételas como a,b, etc. \citep{Bak:63a,Bak:63b}. Por favor,
fíjese que todas las referencias \citep{Garcia:2007}
listadas en este apartado \citep{Garcia:2008}
deben ser citadas directamente en el cuerpo del texto \citep{Garcia:2004}, \citep{Dog:58}, \citep{Keo:58}.

Por favor, tenga en cuenta que las referencias al final de este
documento cumplen con el estilo anteriormente mencionado. Los
artículos que no hayan sido publicados deben ser citados como ``no
publicado.'' Ponga en mayúscula únicamente la primera palabra del
título, excepto el caso de nombres propios y símbolos de elementos.

Si está utilizando LaTeX, puede procesar una base de datos de bibliografía externa o insertarla directamente en la sección de referencias.
Las notas al pie de página se deben evitar en la medida de lo posible.

\subsection{Abreviaciones y acrónimos}

Defina las abreviaciones y acrónimos la primera vez que se usan en
el texto, incluso después de ya hayan sido definidos en el resumen.
Abreviaciones tales como IFAC, SI, ac, y dc no necesitan ser
definidas. Abreviaciones que incorporen periodos no deben tener
espacios: escriba ``C.N.R.S.,'' no ``C. N. R. S.'' No utilice
abreviaciones en el título salvo que sea inevitable (por ejemplo,
``RIAI'' en el título de este artículo).


\begin{table*}[t]
\centering
\caption{Comparación de las especificaciones para cada diseño del sistema. }
\label{tabladeseables}
\begin{tabular}{lcccccc}   \hline
Controlador  & $Re(\lambda)_{max}$  & $u_{max}$ & $t_{est}^{max}$ & $noise_{max}$ & $u_{nom}$ & $t_{est}^{nom}$  \\ \hline
B23 & INA & INA & INA & INA & AD & AIND \\
M23 & AD & AD & AD & T & AD & AIND \\
PPGA23 & \textbf{AD} & \textbf{AD}& \textbf{AD}&\textbf{AD} &\textbf{AD} & \textbf{AD}\\
 \hline
 W34 & AD & AD & D & T & AD & IND \\
 M34 & AD & AD & D & AD & AD & AD \\
 \textbf{PPGA23} & \textbf{AD} & \textbf{AD} & \textbf{AD} & \textbf{AD} & \textbf{AD} & \textbf{AD} \\
 \textbf{PPGA34} & \textbf{AD} & \textbf{AD} & \textbf{AD} & \textbf{AD} & \textbf{AD} & \textbf{AD} \\
  \hline
J45 & AD & IND & AD & IND & AD & AD \\
M45 & AD & AD & IND & T & AD & IND \\
\textbf{PPGA23} & \textbf{D} & \textbf{AD} & \textbf{D} & \textbf{T} & \textbf{AD} & \textbf{D} \\
\textbf{PPGA34} & \textbf{AD} & \textbf{AD} & \textbf{D} & \textbf{D} & \textbf{AD} & \textbf{D}  \\
\textbf{PPGA45} & \textbf{AD} & \textbf{AD} & \textbf{AD} & \textbf{AD} & \textbf{AD} & \textbf{D} \\
 \hline
\end{tabular}
\end{table*}



\subsection{Ecuaciones}

Preste atención a la fuente de las ecuaciones. Si está usando MS-Word, la fuente para las ecuaciones debe ser cambiada a \emph{Texto normal} en el menú \emph{Herramientas de ecuación}.

Numere las ecuaciones consecutivamente con números de ecuaciones
entre paréntesis justificado al margen derecho, como en \eqref{e2}.
Primero use el editor de ecuaciones para crear la ecuación. Después
seleccione el estilo ``Equation''. Presione la tecla de tabulador y
escriba el número de ecuación entre paréntesis. Para hacer sus
ecuaciones más compactas, puede usar el solidus ( / ), la función
exp, o los exponentes apropiados. Utilice paréntesis para evitar
ambigüedades en los denominadores. Ponga signos de puntuación en las
ecuaciones cuando formen parte de una frase, como en este ejemplo
$\int_0^{r_2} F (r, \varphi ) dr d\varphi = [\sigma r_2 / (2 \mu_0 )]$.

Si es una fórmula que no forma parte de un texto, debe numerarse siguiendo el formato:

\begin{equation} \label{e2}
\int_0^{\inf} exp(-\lambda |z_j - z_i |) \lambda^{-1} J_1 (\lambda  r_2 ) J_0 (\lambda r_i ) d\lambda
\end{equation}

Asegúrese de que los símbolos de su ecuación han sido definidos
antes de que la ecuación aparezca o inmediatamente después. Ponga en
cursiva los símbolos (T podría referirse a la temperatura, pero T es
la unidad tesla). Refiérase a ``(1),'' no ``Ec. (1)'' o ``ecuación
(1),'' excepto al principio de la frase: ``la ecuación (1) es ...
.''

\subsection{Otras recomendaciones}

Utilice un espacio tras los periodos y dos puntos. Evite utilizar
participios, tales como, ``Utilizando (1), se calculó el
potencial.'' [No está claro quien o qué usó (1).] En su lugar
escriba ``El Potencial fue calculado empleando (1),'' o ``Empleando
(1), se calculó el potencial.''

\section{Conclusiones}

Una sección de conclusiones no es necesaria. Sin embargo, las conclusiones pueden revisar los puntos más importantes de un artículo, pero no debe replicarse el resumen en las conclusiones. Las conclusiones pueden tratar sobre la importancia del trabajo realizado o sugerir aplicaciones o trabajos futuros.

Repetido. Una sección de conclusiones no es necesaria. Sin embargo, las conclusiones pueden revisar los puntos más importantes de un artículo, pero no debe replicarse el resumen en las conclusiones. Las conclusiones pueden tratar sobre la importancia del trabajo realizado o sugerir aplicaciones o trabajos futuros.

Repetido. Una sección de conclusiones no es necesaria. Sin embargo, las conclusiones pueden revisar los puntos más importantes de un artículo, pero no debe replicarse el resumen en las conclusiones. Las conclusiones pueden tratar sobre la importancia del trabajo realizado o sugerir aplicaciones o trabajos futuros.

%%%%%%%%%%%%%%%%%%%%%%%%%%%%%%%%%%%%%%%%%%%%%%%%%%%%%%%%%%%%%%%%%%%%%%%%%%%%%%%%%%%%%%%%%%

\section*{Agradecimientos}

Este trabajo ha sido realizado parcialmente gracias al apoyo de la Agencia Nacional (los agradecimientos de financiación y apoyos han de ser incluidos aquí).

\label{}

%% The Appendices part is started with the command \appendix;
%% appendix sections are then done as normal sections
%% \appendix

%% \section{}
%% \label{}

%% References
%%
%% Following citation commands can be used in the body text:
%% Usage of \cite is as follows:
%%   \cite{key}         ==>>  [#]
%%   \cite[chap. 2]{key} ==>> [#, chap. 2]
%%


%% References with BibTeX database:

\bibliographystyle{elsarticle-harv}
\bibliography{sci2022}

%% Authors are advised to use a BibTeX database file for their reference list.
%% The provided style file elsarticle-num.bst formats references in the required Procedia style

%% For references without a BibTeX database:

% \begin{thebibliography}{00}

%% \bibitem must have the following form:
%%   \bibitem{key}...
%%

% \bibitem{}

% \end{thebibliography}

\appendix
\section{Primer apéndice}    % Capa Apéndice debe tener un título corto.
Este texto está repetido. Si utiliza Word, use o bien Microsoft Editor de Ecuaciones o
MathType  para las ecuaciones de su artículo (Insertar | Objeto |
Crear Nuevo | Microsoft Editor de Ecuaciones o Ecuación MathType).
No debe seleccionar la opción ``Flotar'' sobre el texto. Por
supuesto, LaTeX gestiona las ecuaciones a través de macros
pre-programadas.

\section{Segundo apéndice}

Este texto está repetido. Use el Sistema Internacional como unidades primarias. Se pueden usar
otras unidades como unidades secundarias (entre paréntesis). Esto se
aplica a artículos sobre almacenamiento de datos. Por ejemplo,
escriba ``$15 Gb/cm^2$'' ($100 Gb/in^2$). Se considera una excepción
cuando las unidades inglesas se usan como identificadores
comerciales, como unidad de disco de 3.5 pulgadas. Evite mezclar
unidades del Sistema Internacional con el Sistema Cegesimal, tales
como corriente en amperios y campo magnético en  oersteds. Esto a
menudo lleva a confusión porque las ecuaciones no son
dimensionalmente equiparables. Si debe usar unidades mezcladas,
especifique claramente las unidades para cada cantidad  en la
ecuación.

La unidad en el Sistema Internacional para la fuerza del campo
magnético H es A/m. Sin embargo, si desea utilizar unidades de $T$,
o bien refiérase a densidad de flujo magnético $B$ o fuerza del
campo magnético simbolizado como $\mu_0 H$. Utilice el punto
centrado para separar unidades compuestas,  es decir, $A\cdot m^2$.

\section{Tercer apéndice}

\subsection{Más sobre figuras y tablas}

Las etiquetas de los ejes de las figuras son a menudo fuentes de
confusión. Utilice palabras en lugar de símbolos. Como ejemplo,
escriba la cantidad ``Magnetización,'' o ``Magnetización M,'' no
sólo ``M.'' Ponga las unidades entre paréntesis. No etiquete los
ejes únicamente con unidades. Por ejemplo,
escriba ``Magnetización (A/m)'' o ``Magnetización (A $\cdot$
m$^{-1}$),'' no sólo ``A/m'' No etiquete los ejes con una relación
de cantidades y unidades. Por ejemplo, escriba ``Temperatura (K),''
no ``Temperatura/K.''

\subsection{Otro subapéndice}

Los multiplicadores pueden ser especialmente fuente de confusión.
Escriba ``Magnetización (kA/m)'' o ``Magnetización ($10^3$ A/m).'' No
escriba ``Magnetización (A/m) $\times$ 1000'' porque el lector no
sabría si la etiqueta es 16000 A/m o
0.016 A/m. Las etiquetas de las figuras deben ser legibles,
aproximadamente de 8 a 12 puntos.

\end{multicols}

\end{document}

%%
%% End of file `ejemplo latex RIAI.tex'.
